\documentclass[10pt,a4paper]{amsart}
\usepackage[margin=19mm]{geometry}
\usepackage{amsmath,amssymb,mathtools}
\usepackage[scr=rsfs,cal=euler]{mathalfa}
\usepackage{xfrac}
\usepackage[unicode,hidelinks,bookmarks=false]{hyperref}
\newcommand{\ie}{i.e.\ }
\newcommand{\cf}{cf.\ }
\newcommand{\resp}{respectively\ }
\newcommand{\Subsection}{\subsection}
\providecommand{\nobreakdash}{\nobreak\hbox{-}}
\setlength{\emergencystretch}{2em}
\multlinegap0pt
\usepackage{bm}
\usepackage{bbm}
\usepackage{dsfont}
\usepackage{xspace}
\usepackage{cancel}
\usepackage{mathtools}
\usepackage{amsthm}
\makeatletter
\@ifundefined{theorem}{\newtheorem{theorem}{Theorem}}{}
\makeatother
\newcommand{\R}{\mathbb{R}}
\title[Chirp-Induced Non-Separable Gabor Windows on $\mathbb{R}^d$]{Chirp-Induced Non-Separable Gabor Windows on $\mathbb{R}^d$}
\author{Lorenzo De Leonardis, Alessandro Mazzoccoli, Pierluigi Vellucci}
\date{Source: arXiv:2606.09300v2 (CC BY 4.0)}
\begin{document}
\maketitle

\noindent\emph{Source and licence.} Chirp-Induced Non-Separable Gabor Windows on $\mathbb{R}^d$. Version arXiv:2606.09300v2 (CC BY 4.0), distributed under the Creative Commons Attribution 4.0 International licence, as recorded in the arXiv licence metadata for this exact version and shown on the abstract page https://arxiv.org/abs/2606.09300v2. Full rights notice: licenses/arxiv-2606.09300-CC-BY-4.0.txt.\par\smallskip

\noindent\textit{Editorial scope.} Counted unit: the Theorem whose source label is \texttt{thm:four\_point\_obstruction} in the article (source0.tex lines 1064--1144, source label \texttt{thm:four\_point\_obstruction}), delivered together with the article's own complete original proof of it (source0.tex lines 1145--1315, one \texttt{proof} environment of 171 lines). No mathematics is written, repaired or re-derived: statement and proof are the article's own text line for line. Required context retained uncounted: none. Every definition, display and label of those lines is retained unchanged. Omitted from this package and not redistributed: 1--1063, 1316--2348. Nothing inside a retained excerpt is omitted or altered: every retained body is delivered exactly as the article's lines are written, so no omission receipt of any kind accompanies this package. Citations: the retained excerpts carry no citation command at all, so no bibliography entry of the article is transcribed and no reference is redistributed. Reference adaptation: none was needed and none was made; every reference inside the delivered unit resolves inside it, so no unresolved reference or citation is produced. Printed slips kept literal: no printed word, symbol, hypothesis or number of the counted statement or of its proof was altered. Layout and class: the article's own class is \texttt{elsarticle} with packages including amsmath, amsthm, amssymb, mathtools, bm, microtype, enumitem, hyperref, xcolor; the wrapper is the lane's pinned \texttt{amsart} editorial wrapper at 10pt on A4, which restates the theorem-like environments the retained text needs and adds the article's own macro definitions that the retained text needs (\texttt{\textbackslash R}), with extra wrapper packages (bm, bbm, dsfont, xspace, cancel, mathtools). The delivered numbering is the unit's own. Assets: no retained span contains a figure environment, a graphics-inclusion command, a PDF-page inclusion, a table or a TikZ picture; no image file is copied into this package, the package declares no asset dependency and no image file is redistributed with it. Licence: version arXiv:2606.09300v2 of this article is distributed under the Creative Commons Attribution 4.0 International licence, as recorded in the arXiv licence metadata for this exact version and shown on the abstract page https://arxiv.org/abs/2606.09300v2. A full rights notice is delivered at licenses/arxiv-2606.09300-CC-BY-4.0.txt. Characters: every retained excerpt is pure ASCII, so the delivered engine, XeLaTeX with its default Latin Modern text font, reports no missing character.
\begin{theorem}[Exact four-point obstruction generated by off-diagonal chirps]
\label{thm:four_point_obstruction}
{Let $d\ge 2$, let $C=(c_{mn})_{m,n=1}^d$ be symmetric, and let
$$
g_C(x)
=
e^{\pi i x^T Cx}\prod_{m=1}^d g_m(x_m),
\qquad x\in\R^d.
$$
Fix two distinct indices $j\neq k$. Assume that, for each $m$, the factor $g_m$
is nowhere zero on a nonempty open interval $I_m\subset\R$. Fix points
$\xi_m\in I_m$ for all $m\neq j,k$, and define the two-variable restriction
$$
F_\xi(u,v)
:=
g_C(x_\xi(u,v)),
$$
where $x_\xi(u,v)\in\R^d$ is given by
$$
(x_\xi(u,v))_j=u,\qquad
(x_\xi(u,v))_k=v,\qquad
(x_\xi(u,v))_m=\xi_m\quad (m\neq j,k).
$$
Then $F_\xi$ is nowhere zero on $I_j\times I_k$, and for every
$u_1,u_2\in I_j$, $v_1,v_2\in I_k$, one has
$$
\frac{
F_\xi(u_1,v_1)F_\xi(u_2,v_2)
}{
F_\xi(u_1,v_2)F_\xi(u_2,v_1)
}
=
e^{2\pi i\,c_{jk}(u_1-u_2)(v_1-v_2)}.
$$
In particular, the four-point quotient is identically equal to $1$ for every
separable function, whereas for the chirped tensor-product window it detects exactly
the off-diagonal coefficient $c_{jk}$.}

{Consequently, for every pair of bounded intervals
$J_j\subset I_j$, $J_k\subset I_k$ of positive length, the normalized rectangular defect
$$
\mathcal R_{jk}(g_C;J_j,J_k)^2
:=
\frac{1}{|J_j|^2|J_k|^2}
\int_{J_j^2}\int_{J_k^2}
\left|
\frac{
F_\xi(u_1,v_1)F_\xi(u_2,v_2)
}{
F_\xi(u_1,v_2)F_\xi(u_2,v_1)
}
-1
\right|^2
\,dv_1\,dv_2\,du_1\,du_2
$$
satisfies the exact formula
$$
\mathcal R_{jk}(g_C;J_j,J_k)^2
=
\frac{4}{|J_j|^2|J_k|^2}
\int_{J_j^2}\int_{J_k^2}
\sin^2\!\left(
\pi c_{jk}(u_1-u_2)(v_1-v_2)
\right)
\,dv_1\,dv_2\,du_1\,du_2.
$$
Hence
$$
\mathcal R_{jk}(g_C;J_j,J_k)=0
\quad\Longleftrightarrow\quad
c_{jk}=0.
$$
Moreover, if $L_j:=|J_j|$ and $L_k:=|J_k|$, then, as $c_{jk}\to0$,
$$
\mathcal R_{jk}(g_C;J_j,J_k)^2
=
\frac{\pi^2}{9}\,c_{jk}^2 L_j^2L_k^2
+
O\!\left(c_{jk}^4 L_j^4L_k^4\right).
$$}
\end{theorem}

\begin{proof}
{Since each $g_m$ is nowhere zero on $I_m$, the restricted function
$F_\xi$ is nowhere zero on $I_j\times I_k$. We now compute the four-point quotient.}

{Write $x=x_\xi(u,v)$. The quadratic form $x^TCx$, restricted to the variables
$(u,v)$, can be decomposed as
$$
x_\xi(u,v)^T C x_\xi(u,v)
=
c_{jj}u^2+c_{kk}v^2+2c_{jk}uv+\alpha u+\beta v+\rho,
$$
where $\alpha,\beta,\rho$ depend on $C$ and on the fixed coordinates
$\xi_m$, but not on $u$ and $v$ in any other way. Hence
$$
F_\xi(u,v)
=
e^{\pi i(c_{jj}u^2+c_{kk}v^2+2c_{jk}uv+\alpha u+\beta v+\rho)}
\,g_j(u)g_k(v)
\prod_{m\neq j,k}g_m(\xi_m).
$$
Consider now
$$
Q(u_1,u_2,v_1,v_2)
:=
\frac{
F_\xi(u_1,v_1)F_\xi(u_2,v_2)
}{
F_\xi(u_1,v_2)F_\xi(u_2,v_1)
}.
$$
All one-variable factors cancel:
the factors involving $g_j$, $g_k$, the fixed coordinates, the terms
$c_{jj}u^2$, $c_{kk}v^2$, $\alpha u$, $\beta v$, and $\rho$
appear equally in numerator and denominator. The only surviving term is the mixed
term $2c_{jk}uv$. Therefore
$$
\begin{aligned}
Q(u_1,u_2,v_1,v_2)
&=
e^{
\pi i\,2c_{jk}
\big(
u_1v_1+u_2v_2-u_1v_2-u_2v_1
\big)
}\\
&=
e^{
2\pi i\,c_{jk}(u_1-u_2)(v_1-v_2)
}.
\end{aligned}
$$
This proves the first assertion.}

{We also recall the elementary reason why this quotient is an obstruction to
separability. Let $H$ be any nowhere-zero separable function on a rectangle
$J_j\times J_k$, say
$$
H(u,v)=a(u)b(v),
\qquad u\in J_j,\ v\in J_k,
$$
for some one-variable functions $a$ and $b$. Then, for all
$u_1,u_2\in J_j$ and $v_1,v_2\in J_k$,
$$
\frac{
H(u_1,v_1)H(u_2,v_2)
}{
H(u_1,v_2)H(u_2,v_1)
}
=
\frac{
a(u_1)b(v_1)a(u_2)b(v_2)
}{
a(u_1)b(v_2)a(u_2)b(v_1)
}
=
1.
$$
Thus the quantity
$$
\frac{
H(u_1,v_1)H(u_2,v_2)
}{
H(u_1,v_2)H(u_2,v_1)
}
-1
$$
vanishes identically for separable functions. Applied to $H=F_\xi$, the formula
computed above shows that the obstruction is generated exactly by the mixed chirp
coefficient $c_{jk}$.}

{Using the identity
$$
|e^{i\theta}-1|^2=4\sin^2(\theta/2),
$$
with
$$
\theta=2\pi c_{jk}(u_1-u_2)(v_1-v_2),
$$
we obtain the exact formula
$$
\mathcal R_{jk}(g_C;J_j,J_k)^2
=
\frac{4}{|J_j|^2|J_k|^2}
\int_{J_j^2}\int_{J_k^2}
\sin^2\!\left(
\pi c_{jk}(u_1-u_2)(v_1-v_2)
\right)
\,dv_1\,dv_2\,du_1\,du_2.
$$
This quantity is zero if $c_{jk}=0$. Conversely, suppose it is zero. Since the integrand is
continuous and nonnegative, it must vanish identically on
$J_j^2\times J_k^2$. Hence
$$
\sin\!\left(
\pi c_{jk}(u_1-u_2)(v_1-v_2)
\right)=0
$$
for all $u_1,u_2\in J_j$ and $v_1,v_2\in J_k$. Because $J_j$ and $J_k$
have positive length, the product
$(u_1-u_2)(v_1-v_2)$ ranges over a nontrivial interval containing $0$.
The only way the above sine can vanish identically on such a continuum is $c_{jk}=0$.
Thus
$$
\mathcal R_{jk}(g_C;J_j,J_k)=0
\quad\Longleftrightarrow\quad
c_{jk}=0.
$$
Finally, for $c_{jk}\to0$, we use
$$
\sin^2 z=z^2+O(z^4).
$$
Therefore
$$
\begin{aligned}
\mathcal R_{jk}(g_C;J_j,J_k)^2
&=
\frac{4\pi^2c_{jk}^2}{|J_j|^2|J_k|^2}
\left(
\int_{J_j^2}(u_1-u_2)^2\,du_1\,du_2
\right)
\left(
\int_{J_k^2}(v_1-v_2)^2\,dv_1\,dv_2
\right)\\
&\qquad
+
O\!\left(c_{jk}^4 L_j^4L_k^4\right).
\end{aligned}
$$
For an interval $J$ of length $L$,
$$
\int_{J^2}(s-t)^2\,ds\,dt=\frac{L^4}{6}.
$$
Hence
$$
\mathcal R_{jk}(g_C;J_j,J_k)^2
=
\frac{4\pi^2c_{jk}^2}{L_j^2L_k^2}
\frac{L_j^4}{6}\frac{L_k^4}{6}
+
O\!\left(c_{jk}^4 L_j^4L_k^4\right),
$$
that is,
$$
\mathcal R_{jk}(g_C;J_j,J_k)^2
=
\frac{\pi^2}{9}\,c_{jk}^2L_j^2L_k^2
+
O\!\left(c_{jk}^4L_j^4L_k^4\right).
$$
This completes the proof.}
\end{proof}

\end{document}
